% concatenated from main.tex
\documentclass[11pt,a4paper,longbibliography]{amsart}
\usepackage[margin=2.5cm]{geometry}
\usepackage{cite}
\usepackage{amsmath,amsfonts,amssymb,amsthm,amscd,mathrsfs,mathdots,bbm,color,bm,enumerate,comment,graphicx,newtxmath,newtxtext,subcaption,stmaryrd,caption,soul, xcolor}
\usepackage{comment}
\usepackage{physics}
\usepackage{multirow}
\usepackage[normalem]{ulem}

\numberwithin{equation}{section}
\pagenumbering{arabic} 

\definecolor{equationcolor}{RGB}{222,94,100}
\definecolor{alecolor}{RGB}{198,113,190}


\usepackage[linecolor=white,backgroundcolor=orange,bordercolor=red]{todonotes}

\usepackage[colorlinks=true,linkcolor=equationcolor,citecolor=teal]{hyperref}

\hypersetup{urlcolor=teal}
\definecolor{equationcolor}{RGB}{222,94,100}

% \DeclareFontFamily{U}{mathx}{}
% \DeclareFontShape{U}{mathx}{m}{n}{<-> mathx10}{}
% \DeclareSymbolFont{mathx}{U}{mathx}{m}{n}
% \DeclareMathAccent{\widecheck}{0}{mathx}{"71}

\parskip3pt

\def\todon{\todo[noline]}
\def\todoq{\todo[backgroundcolor=cyan,bordercolor=green,noline]}    

\def\R{{\mathbb R}}
\def\C{{\mathbb C}}
\def\N{{\mathbb N}}
\def\Z{{\mathbb Z}}
\def\UT{\mathrm{UT}}

\def\oprod{\overset{\leftarrow}{\prod} }
\def\End{{\mathrm End}}

\def\Sym{{\mathrm Sym}}

\def\ts{\otimes}

\def\tk{\tilde{k}}

\def\EE{{\mathbb E}}
\def\PP{{\mathbb P}}
\def\M{{\mathbb M}}
\def\AA{\mathcal{A}}
\def\BB{\mathcal{B}}
\def\FF{\mathcal{F}}
\def\SS{\mathcal{S}}

\def\de{ \mathrm{d}}
\def\l{\lambda}
\def\L{\Lambda}
\def\id{\mathrm{Id}}
\def\Tr{\mathrm{Tr}}
\def\tr{\mathrm{tr}}
\def\Wg{\mathrm{Wg}}
\def\sinc{\operatorname{sinc}}

\def\Ton{\mathrm{T}}

\def\Hm{\mathrm{H}^{\uparrow}}
\def\Wm{\mathcal{W}^{\uparrow }}
\def\Fm{\mathcal{F}^{\uparrow}}
\def\l{\lambda}
\def\b{\beta}
\def\Ch{\operatorname{Ch}}
\def\Meix{\operatorname{Meix}}
\def\Kr{\operatorname{Kr}}
\def\rhosc{\rho_{\mathrm{sc}}}
\def\sine{\operatorname{sine}}
\def\Ai{\operatorname{Ai}}

\def\sigmaPc{{\sigma}_{P_{<N}}^{\hbar}}
\def\sigmaHc{{\sigma}_{H_N}^{\hbar}}
\def\sigmaKc{{\sigma}_{P_{<N}}^{\hbar}}
\def\sigmaQc{{\sigma}_{H_N}^{\hbar}}
\def\psigmaKc{\widecheck{p\sigma_{K_N}}}
\def\chiAi{\chi^{\Ai_1,N}_D}

\def\Hsm{\mathrm{H}^{\uparrow \hspace{-0,05 cm}\uparrow}}
\def\Wsm{\mathcal{W}^{\uparrow \hspace{-0,05 cm}\uparrow}}
\def\Fsm{\mathcal{F}^{\uparrow \hspace{-0,05 cm}\uparrow}}
\def\conv{\operatorname{conv}}
\def\ch{\operatorname{ch}}

\def\R{{\mathbb R}}
\def\Z{{\mathbb Z}}
\def\N{{\mathbb N}}
\def\C{{\mathbb C}}
\def\rme{{\mathrm e}}
\def\e{\mathbbm{e}}
\def\UT{{\mathrm{UT}}}
\def\Dir{\mathrm{Dir}}
\def\rcat{\xrightarrow[\mathrm{cat}]{}}


\def\l{\lambda}
\def\b{\lambda}
\def\um{\overset{UT}{\succ}}

\renewcommand{\P}{\mathbb{P}}
\newcommand{\E}{\mathbb{E}}
\newcommand{\T}{\mathbb{T}}

\def\eqlaw{\overset{(\mathrm{law})}{=}}

\usepackage{thmtools,amsthm}
\theoremstyle{plain}

\newtheorem{thm}{Theorem}[section]
\newtheorem{conj}{Conjecture}
\newtheorem{cor}[thm]{Corollary}
\newtheorem{prop}[thm]{Proposition}
\newtheorem{lem}[thm]{Lemma}
\newtheorem{defn}{Definition}[section]
\newtheorem{ex}{Example}[section]

\newtheorem{mainthm}{Theorem}
\renewcommand{\themainthm}{\Alph{mainthm}}

\theoremstyle{definition}
\newtheorem*{notation}{Notation}
\theoremstyle{remark}
\newtheorem{rem}{Remark}
\newtheorem{fact}{Fact}
% \newtheorem{mrem}{\textbf{\rev{Main Result}}}


\declaretheorem[
name=Theorem,
style=bold,
numberwithin=section,
refname={Theorem,A},
Refname={Theorem,A}
]{maintheorem}

\newtheorem*{theoremA}{Theorem A}
\newtheorem*{theoremB}{Theorem B}

\newcommand{\rev}[1]{{\color{red}#1}}
\newcommand{\REV}[1]{{\textbf{\color{red}[[#1]]}}}
\newcommand{\revG}[2]{{\color{red}\st{#1}} {\color{blue} #2}}


\newcommand{\cy}[1]{{\color{blue!50!cyan}#1}}
\newcommand{\jcz}[1]{\cy{(\texttt{JCz:} \textit{#1})}}
\newcommand{\gr}[1]{{\color{green!50!black}(\texttt{GG:} \textit{#1})}}
\newcommand{\fdc}[1]{{\color{red}(\texttt{FDC:} \textit{#1})}}
\newcommand{\ml}[1]{{\color{orange}(\texttt{ML:} \textit{#1})}}
\newcommand{\rd}[1]{{\color{red}#1}}



\newcommand{\blue}[1]{{\color{blue}#1}}

% \newcommand{\be}{\begin{equation}}
	% \newcommand{\ee}{\end{equation}}

% \newcommand{\ben}{\begin{equation*}}
	% \newcommand{\een}{\end{equation*}}

\def\_#1{\def\next{#1}
	\ifx\next\risingsign\expandafter\rising\else^{\underline{#1}}\fi}
\def\risingsign{^}
\def\rising#1{^{\overline{#1}}}

% \title[Distribution-Preserving Stochastic Semigroups]{When Limits Do Not Commute: Distribution-Preserving Stochastic Semigroups \fdc{Title with `limit', `distribution-preserving stochastic semigroups' and `block-triangular majorisation'. }}
\title[Limits of  Stochastic Semigroups and Block-Triangular Majorisation]{Limits of Stochastic Semigroups and Block-Triangular Majorisation}

\author[F.~D.~Cunden]{Fabio Deelan Cunden}
\address{Dipartimento di Matematica, Universit\'a degli Studi di Bari, I-70125 Bari, Italy, and
	INFN, Sezione di Bari, I-70126 Bari, Italy}
\email{fabio.cunden@uniba.it}

\author[J. Czartowski]{Jakub Czartowski}
\address{School of Physics, Trinity College Dublin, Dublin 2, Ireland}
\address{School of Physical and Mathematical Sciences, Nanyang Technological University, 21 Nanyang Link, 637371 Singapore, Republic of Singapore}
\email{jakub.czartowski@tcd.ie}

\author[G.~Gramegna]{Giovanni Gramegna}
\address{Dipartimento di Fisica, Universit\`a degli Studi di Bari, I-70126 Bari, Italy, and
	INFN, Sezione di Bari, I-70126 Bari, Italy}
\email{giovanni.gramegna@uniba.it}

\author[M. Ligab\`o]{Marilena Ligab\`o}
\address{Dipartimento di Matematica, Universit\'a degli Studi di Bari, I-70125 Bari, Italy}
\email{marilena.ligabo@uniba.it}
\thanks{This research is supported by Gruppo Nazionale di Fisica Matematica GNFM-INdAM, Istituto Nazionale di Fisica Nucleare INFN through the project QUANTUM, and financially supported by PNRR MUR project PE0000023-NQSTI, PNRR MUR project CN00000013 ‘Italian National Centre on HPC, Big Data
	and Quantum Computing’, PRIN 2022 project 2022TEB52W-PE1-
	`The charm of integrability: from nonlinear waves to random matrices', the project ``Patto territoriale sistema universitario pugliese'', the start-up grant associated with the Nanyang Assistant Professorship awarded to Nelly Ng at Nanyang Technological University, Singapore, and Taighde Éireann -- Research Ireland under Grant No.~IRCLA/2022/3922.}

\allowdisplaybreaks

\begin{document}
	
	\begin{abstract}  
		We investigate limits of semigroups of stochastic matrices defined by their invariant distribution. Given probability vectors $\gamma(\beta)$ depending on a parameter $\beta$, we introduce a notion of convergence as $\beta\to\infty$ for the corresponding semigroups of $\gamma(\beta)$-preserving stochastic matrices and investigate the structure of the resulting limit. In general, the limiting semigroup differs from the semigroup preserving the limiting distribution, showing that these two operations do not commute.
		We develop a general framework for such limiting semigroups and study in detail the case in which the invariant distributions are Gibbs vectors at the inverse temperature $\beta$. We show that the limiting semigroup consists of block-upper-triangular stochastic matrices subject to additional substochasticity constraints. We characterise and enumerate their extremal elements and determine the preorder on probability vectors induced by the action of the semigroup.
		The resulting notion of Block-Triangular majorisation interpolates between ordinary majorisation and upper triangular (aka unordered) majorisation. We show that it is completely characterised by a finite family of monotones and analyse the corresponding behaviour of R{\'e}nyi $\alpha$-entropies as $\beta\to\infty$.
	\end{abstract}
	
	\maketitle
	
	
	\section{Introduction}
	\label{sec:intro}
	Majorisation governs a number of fundamental state-conversion problems in quantum information theory and statistical physics. Examples include bipartite pure-state entanglement transformations~\cite{nielsen1999conditions}, coherence manipulation~\cite{baumgratz2014quantifying}, and thermodynamic state transitions under Gibbs-preserving constraints~\cite{horodecki2013fundamental}. Among these, the resource theory of athermality exhibits the richest majorisation structure, described by thermomajorisation~\cite{gour2015resource,deoliveirajunior2022}.
	
	While much recent work has focused on coherent systems, especially in asymptotic \cite{shiraishi2025recovery} and catalytic settings \cite{son2024hierarchy}, the energy-incoherent regime continues to reveal nontrivial structural aspects of finite-size thermodynamics. In particular, degeneracies in the energy spectrum introduce additional geometric and combinatorial features that are not present in the nondegenerate case.
	
	For a given nonincreasing sequence of energy levels $E_1\geq E_2\geq\cdots \geq E_n$, consider the associated Gibbs probability vector at each inverse temperature $\beta\geq0$,
	\begin{equation}
		\label{eq:Gibbs_intro}
		\gamma(\beta)
		=
		\frac1{Z(\beta)}
		\bigl(
		e^{-\beta E_1},
		\ldots,
		e^{-\beta E_n}
		\bigr),
		\qquad
		Z(\beta)=\sum_{i} e^{-\beta E_i}.
	\end{equation}
	Let $\mathcal S^{\gamma(\beta)}$ denote the set of $n\times n$ row-stochastic matrices satisfying
	$
	\gamma(\beta)P=\gamma(\beta)
	$.
	Under matrix multiplication, $\mathcal S^{\gamma(\beta)}$ is a semigroup with identity. Equivalently, it is the semigroup of Markov transition matrices having $\gamma(\beta)$ as a stationary distribution. The convex geometry of this set, its extreme points and its relation to other notions such as transportation polytopes has been extensively studied in the literature~\cite{Hartfiel1974, GregoryKirklandPullman1992, vom2019d,Cunden2025SemigroupMajorisation}.
	
	A natural question is: what becomes of the finite-temperature semigroups
	$\mathcal{S}^{\gamma(\beta)}$
	in the zero-temperature limit, $\beta\to\infty$? 
	As the temperature decreases, the Gibbs distribution $\gamma(\beta)$ becomes uniformly supported in the lowest energy space. Let $\gamma(\infty)$ be this limit. If we denote by $S^{\gamma(\infty)} $ the set of $\gamma(\infty)$-preserving row-stochastic matrices, is it true that
	\begin{equation}
		\label{eq:question}
		\lim_{\beta\to\infty}S^{\gamma(\beta)}=S^{\gamma(\infty)}?
	\end{equation}
	At low temperatures, the Gibbs state becomes increasingly concentrated on the ground-state sector. Since absolute zero is an ideal limit, it is natural to study the behaviour of the semigroups $\mathcal S^{\gamma(\beta)}$ as $\beta\to\infty$. Equation~\eqref{eq:question} asks whether taking the zero-temperature limit commutes with passing from a Gibbs state to the semigroup of stochastic maps that preserve its limit $\gamma(\infty)$.
	
	In this paper we show that the answer is negative: taking the limit of invariant distributions and taking the semigroup of stochastic matrices preserving them are, in general, non-commuting operations. Different families of invariant (not necessarily Gibbs) distribution $\gamma(\beta)$ may converge to the same limiting distribution while giving rise to different limiting semigroups. Understanding precisely which features of the asymptotic behaviour determine the limit, and which are irrelevant, is the driving question of the present work. In this paper, we first clarify the precise meaning of the limit on the left-hand side of~\eqref{eq:question}, that is, the limit of a family of semigroups of matrices indexed by the parameter $\beta\ge0$, and then compute the limit.
	
	\subsection{Outline of the paper}
	\label{sec:outline}
	The paper is organised as follows. We start in Section~\ref{sec:preparatory} by setting the notation, terminology and definitions on stochastic matrices, semigroup majorisation, and Kuratowski convergence for sequences of sets. In Section~\ref{sec:results} we present the results. In Section~\ref{sec:examples} we illustrate the zero temperature limits ($\beta\to\infty$) of Gibbs-preserving semigroups and present the corresponding finite-temperature corrections for low-dimensional explicit examples. The proofs are contained in Section~\ref{sec:proofs}.
	\section{Preliminaries}
	\label{sec:preparatory}
	
	\subsection{Stochastic matrices}
	
	We briefly recall the required notions from majorisation theory; a standard reference is~\cite{marshall1979inequalities}.
	We write $\R_+$ and $\R_{++}$ for the sets of nonnegative and strictly positive real numbers, respectively. Let $\e=(1,\dots,1)\in\R^n$ represent the flat (unnormalised) distribution over $n$ outcomes.
	
	A $n\times n$ matrix $P$ is called (row-)stochastic if $P_{ij}\geq0$ for all $i,j=1,\ldots, n$, and its rows sum to $1$:
	\begin{equation}
		P \e^T=\e^T
	\end{equation}
	Given $q\in\R_+^n$, a stochastic matrix is said to be $q$-preserving if
	\begin{equation}
		qP=q.
	\end{equation}
	
	
	We denote by $\mathcal{S}$ the semigroup of $n\times n$ stochastic matrices and by $\mathcal{S}^q$ the subsemigroup of $q$-preserving stochastic matrices. Both the identity matrix and the rank-one matrix $\e^T q$ (all rows equal to $q$) belong to $\mathcal{S}^q$, so the latter is nonempty. 
	The set $\mathcal{S}^q$ is a compact convex polytope in $\R^{n\times n}$ of dimension at most $(n-1)^2$, with equality for $q\in\R_{++}^n$.
	
	A stochastic matrix $P$ is called {$q$-subpreserving} if $qP\leq q$,
	where the inequality is understood componentwise.
	In the special case $q=\e$, the condition
	$\e P\leq \e$
	is equivalent to requiring that all column sums of $P$ are at most~$1$. In this case, $P$ is said to be {column-substochastic}.
	
	
	\subsection{Semigroup majorisation}
	
	Let
	\begin{equation}
		\Delta_{n-1}
		:=
		\left\{
		x\in\R_+^n :
		\sum_i x_i =1
		\right\}
	\end{equation}
	denote the probability simplex.
	Every stochastic matrix acts on $\Delta_{n-1}$ by right multiplication. Let $\mathcal{T}\subseteq\mathcal{S}$ be a subsemigroup of $\mathcal{S}$.
	
	\begin{defn}[Semigroup majorisation]
		For $x,y\in\Delta_{n-1}$, we say that $x$ is $\mathcal{T}$-majorised by $y$, and write
		$  x\prec^{\mathcal T} y$,
		if there exists $P\in\mathcal T$ such that $  x=yP$.
	\end{defn}
	
	If $\mathcal T$ is the set of doubly stochastic matrices, one recovers ordinary majorisation:
	\begin{equation}
		x\prec y
		\iff
		\sum_{j=1}^k x_j^\downarrow
		\leq
		\sum_{j=1}^k y_j^\downarrow,
		\qquad
		k=1,\dots,n,
	\end{equation}
	where $x^\downarrow$ denotes the nonincreasing rearrangement of $x$, i.e. $x_1^\downarrow\geq x_2^\downarrow \geq \dots \geq x_n^\downarrow$.
	
	If $\mathcal T$ is the semigroup of upper triangular stochastic matrices, one obtains upper triangular (UT) majorisation, also called unordered majorisation:
	\begin{equation}
		x\prec^{\mathrm{UT}} y
		\iff
		\sum_{j=1}^k x_j
		\leq
		\sum_{j=1}^k y_j,
		\qquad
		k=1,\dots,n.
	\end{equation}
	
	Finally, if $\mathcal T=\mathcal S^q$, we write $    x\prec^q y$.
	When $q$ is a Gibbs probability vector, this relation is known as thermomajorisation.  See~\cite{horodecki2013fundamental, vom2019d} for more details on the structure of $q$-majorisation and its appearance in the context of resource theory of athermality.
	
	\subsection{Kuratowski convergence}
	As we will be concerned with the low-temperature limits of semigroups $\mathcal{S}_n^q$, we recall the appropriate notions of lower and upper set limits following Painlev\'e and Kuratowski~\cite[Chap.~29]{KuratowskiI}.
	
	
	Let $A_1,A_2,\ldots$ be a sequence of subsets of a metric space $(X,d)$.
	\begin{defn}
		The point $p$ belongs to the \textbf{lower limit} $\operatorname{Li} A_k$ of a sequence of sets $A_1,A_2,\ldots$,  if every neighbourhood of $p$ intersects $A_k$ for all but finitely many $k$. 
	\end{defn}
	\begin{defn}
		The point $p$ belongs to the \textbf{upper limit} $\operatorname{Ls} A_k$ of a sequence of sets $A_1,A_2,\ldots$, if every neighbourhood of $p$ intersects $A_k$ for infinitely many $k$. 
	\end{defn}
	
	Equivalently, $   p\in \operatorname{Li} A_k$ if and only if there exist points $p_k\in A_k$ such that $p_k\to p$. 
	Similarly, $p\in \operatorname{Ls} A_k$ 
	if and only if there exists a subsequence $(m_k)$ and points $p_{m_k}\in A_{m_k}$ such that $p_{m_k}\to p$.
	Elements of $\operatorname{Li}A_k$ are called `limit points' of  $(A_k)_{k\geq1}$, and elements of $\operatorname{Ls}A_k$ are called `cluster points' of $(A_k)_{k\geq1}$.
	The lower and upper limits are closed sets, and one always has
	\begin{equation}
		\operatorname{Li} A_k
		\subseteq
		\operatorname{Ls} A_k.
	\end{equation}
	
	\begin{defn}
		The sequence $(A_k)$ is said to be \textbf{convergent} to $A$ in the sense of Kuratowski if $\operatorname{Li} A_k=A =\operatorname{Ls} A_k$. We then write $A=\operatorname{lim}A_k$.
	\end{defn}
	The preceding definitions extend verbatim to nets $(A_i)_{i\in I}$. In this paper, we shall only consider families of sets indexed by positive reals $\beta\geq 0$, and upper and lower limits are always understood in the sequential sense as $\beta\to\infty$. More precisely, for a family $(A_\beta)$, we write $\operatorname{Ls}A_\beta$ for the set $\operatorname{Ls}A_{\beta_k}$ whenever the latter is independent of the choice of sequence $(\beta_k)$ satisfying $\beta_k\to\infty$. The same convention applies to $\operatorname{Li}A_\beta$.
	
	Let $\mathcal F(X)$ denote the set of closed subsets of $X$, the limit being understood in the sense of Kuratowski. The space $\mathcal{F}(X)$ may fail to be topological. However, if $(X,d)$ is compact, then $\mathcal F(X)$ is metrizable via the Hausdorff distance $d_H$ (see, e.g.,~\cite{Burago2001}).
	In this case,
	\begin{equation}
		\lim A_k= A
		\iff
		d_H(A_k,A)\to0.
	\end{equation}
	Assume henceforth that $(X,d)$ is compact.
	\begin{defn}
		A map
		\begin{equation}
			f:(Y,d')\to (\mathcal F(X),d_H)
		\end{equation}
		is said to be \textbf{Painlev\'e--Kuratowski continuous} at $y\in Y$ if
		\begin{equation}
			d'(y_n, y)\to 0
			\quad\Longrightarrow\quad
			d_H\bigl(f(y_n),f(y)\bigr)\to0.
		\end{equation}
		Equivalently,  $y_n\to y$ implies     $f(y_n)\to f(y)$,
		in the sense of Kuratowski.
	\end{defn}
	
	
	\section{Results}
	\label{sec:results} 
	We consider the case of a continuous curve   
	$$
	\gamma \colon \R_+ \to \Delta_{n-1}, \qquad \beta \mapsto \gamma(\beta),
	$$
	converging as $\beta\to\infty$ to a boundary point of $\Delta_{n-1}$.   The probability vectors $\gamma(\beta)$  may not necessarily be Gibbs vectors. For definiteness, we always assume the components of $\gamma(\beta)$ to be nonincreasing. Theorem~\ref{thm:general}  presents the Kuratowski limit as $\beta\to\infty$ of the corresponding  $\gamma(\beta)$-preserving semigroups $\mathcal{S}^{\gamma(\beta)}$. In particular, the limit depends on the multiscale profile of $\gamma(\beta)$ (see Sec.~\ref{sec:multiscale} below), and this shows that the map $\Delta_{n-1}\ni q\mapsto \mathcal{S}^q$ is not continuous at the boundary of $\Delta_{n-1}$. Figure~\ref{fig:cartoon} illustrates the setting of Theorem~\ref{thm:general} in dimension $n=3$.
	
	\begin{figure}
		\centering
		\includegraphics[width=.9\linewidth]{non_cartoon.pdf}
		\caption{ Dependence of the limiting distribution-preserving semigroup on the multiscale profile. Two families of probability vectors, $\gamma(\beta)$ and $\widetilde\gamma(\beta)$, converge to the same boundary point of $\Delta_2$ as $\beta\to\infty$, while approaching it along different asymptotic trajectories. Consequently, their associated semigroups of preserving stochastic matrices have different Kuratowski limits, illustrating that the limiting distribution alone does not determine the limiting semigroup. In particular, limiting semigroup in degenerate case ($x=1$, blue line) is different from non-degenerate cases ($x = 1.5, 2$; magenta and red lines, respectively).
		}
		\label{fig:cartoon}
	\end{figure}
	Theorem~\ref{thm:main} is the specialisation of Theorem~\ref{thm:general}  to the $\beta\to\infty$ limit of $\mathcal{S}^{\gamma(\beta)}$, when $\gamma(\beta)$ is a Gibbs vector with arbitrary multiplicities of the energies. In this case the limit semigroup only depends on the degeneracy structure of $\gamma(\beta)$, and it is a set of block-upper-triangular matrices that we dub BT-stochastic matrices (`BT' for `block triangular'). We  identify and enumerate the extremal BT-stochastic matrices  (Proposition~\ref{prop:extreme_pts}). We then characterise the induced semigroup majorisation in terms of monotones (Theorem~\ref{prop:monotones}), and we study the zero-temperature limit of the R\`enyi entropic inequalities (Proposition~\ref{prop:fate_monotones}). The limiting semigroups and their extreme points are illustrated in Section~\ref{sec:examples} for $n=3$ and $n=4$, followed by a discussion of finite-temperature corrections.	
	
	\subsection{Multiscale profiles and associated limiting stochastic matrices}
	\label{sec:multiscale}
	A \textbf{composition} of $n$ is a sequence of positive integers
	$\mu=(m_1,\ldots,m_k)
	$
	such that $m_1+\cdots+m_k=n$.
	We call $\mu$ a $k$-composition when exactly $k$ parts appear. There are$ 
	\binom{n-1}{k-1}
	$
	such compositions, and altogether $2^{n-1}$ compositions of $n$. See~\cite{Stanley2011}.
	The role of compositions here is to encode the degeneracy structure of  energy vectors or, more generally, to encode the vanishing rates of the components of $\gamma(\beta)$. %It is easier to illustrate this in the cse of Gibbs vectors.
	
	Let $\gamma(\beta)$ be a $n$-probability vector parametrised by $\beta\geq0$. Assume the components of $\gamma(\beta)$ to be nonincreasing. Suppose that there is $k\geq1$ such that $\gamma$ admits a decomposition 
	$$
	\gamma(\beta)
	=
	\bigl(
	\gamma^1(\beta),\ldots,\gamma^k(\beta)
	\bigr)\in\Delta_{n-1},
	$$
	in blocks $\gamma^a(\beta)\in\mathbb{R}_{+}^{m_a}$, $a=1,\ldots,k$ that are separated in scale, in the sense that
	$$
	\|\gamma^1(\beta)\|_1
	\ll
	\|\gamma^2(\beta)\|_1
	\ll\cdots\ll
	\|\gamma^k(\beta)\|_1,
	\qquad
	\beta\to\infty,
	$$
	that is,
	$$
	\frac{\|\gamma^a(\beta)\|_1}
	{\|\gamma^{a+1}(\beta)\|_1}
	\to 0,
	\qquad
	a=1,\ldots,k-1.
	$$
	In this case we say that $\gamma(\beta)$ is of \textbf{type} $\mu=(m_1,\ldots, m_k)$ as $\beta\to\infty$.
	
	We write $\mathcal{S}^{\gamma(\beta)}$
	for the semigroup of row-stochastic $\gamma(\beta)$- preserving matrices. 
	To describe the $\beta\to\infty$ limit of $\mathcal{S}^{\gamma(\beta)}$, it is convenient to organise matrices according to the type of $\gamma(\beta)$.
	We thus partition an $n\times n$ matrix into blocks,
	$$
	P=\bigl(P^{ab}\bigr)_{a,b=1}^k,
	$$
	where the block $P^{ab}$ has size $m_a\times m_b$. Indices inside the blocks will rarely matter and will be suppressed whenever possible. 
	
	
	Suppose furthermore that, for each $a=1,\ldots,k$,
	$$
	\frac{\gamma^a(\beta)}
	{\|\gamma^a(\beta)\|_1}
	\longrightarrow
	q^a,
	$$
	for some probability vector $q^a\in\Delta_{m_a-1}$.
	We call the tuple $(q^1,\ldots,q^k)$ 
	the \textbf{multiscale profile} of the family $\left(\gamma(\beta)\right)_{\beta\geq0}$. As $\beta\to\infty$, 
	\begin{equation}
		\label{eq:limit_gamma_gen}
		\gamma(\beta)
		\to
		\gamma(\infty)
		=
		\big(
		\underbrace{0,\ldots,0}_{n-m_k},
		q^{k}_1,\ldots,q^{k}_{m_k}
		\big).
	\end{equation}
	The semigroup $\mathcal{S}^{\gamma(\infty)}$ preserving the limiting  vector~\eqref{eq:limit_gamma_gen}
	consists of row-stochastic matrices with blocks
	\begin{equation}
		P^{k1}=0, P^{k2}=0,\ldots,  P^{k,{k-1}}=0  
	\end{equation}
	while the lower-right $m_k\times m_k$ block $P^{kk}$ is $q^k$-preserving.
	\begin{ex}
		Given a nonincreasing sequence $
		E=(E_1,\ldots,E_n)\in\mathbb R^n$, with $E_1\geq E_2\geq\cdots \geq E_n$,
		we say that $E$ is of type $\mu=(m_1,\ldots,m_k)$
		if
		\begin{equation}
			\begin{aligned}
				E_1&=\cdots=E_{m_1},\\
				E_{m_1+1}&=\cdots=E_{m_1+m_2},\\
				&\ \vdots\\
				E_{m_1+\cdots+m_{k-1}+1}
				&=\cdots=
				E_{m_1+\cdots+m_k},
			\end{aligned}
		\end{equation}
		with strict inequalities between consecutive blocks, $E_{m_1+\cdots+m_i} > E_{m_1+\cdots+m_i + 1}$. Thus $E$ has exactly $k$ distinct energy levels, with multiplicities $m_1,\ldots,m_k$. 
		Define the \textbf{Gibbs vector} at each inverse temperature $\beta\geq0$,
		\begin{equation}
			\label{eq:Gibbs}
			\gamma(\beta)
			=
			\frac1{Z(\beta)}
			\bigl(
			e^{-\beta E_1},
			\ldots,
			e^{-\beta E_n}
			\bigr),
			\qquad
			Z(\beta)=\sum_{i=1}^n e^{-\beta E_i}.
		\end{equation}
		Let $\mathcal{S}^{\gamma(\beta)}$
		be the semigroup of $\gamma(\beta)$- preserving matrices. 
		To describe the $\beta\to\infty$ limit of $\mathcal{S}^{\gamma(\beta)}$. 
		As before, it is convenient to organise matrices according to the degeneracy pattern $
		\mu=(m_1,\ldots,m_k)$.
		Note that the Gibbs vector decomposes as
		$$
		\gamma(\beta)
		=
		\bigl(
		\gamma^1(\beta),
		\ldots,
		\gamma^k(\beta)
		\bigr),
		$$
		with each block $\gamma^a(\beta)$ of length $m_a$ and equal entries.
		
		Then, for each $a=1,\ldots, k-1$,
		\begin{equation}
			\frac{\|\gamma^a(\beta)\|_1}
			{\|\gamma^{a+1}(\beta)\|_1}=\frac{m_a}{m_{a+1}} \exp\left(-\beta \left(E_{m_1+\cdots+{m_a}}-E_{m_1+\cdots+{m_a}+1}\right)\right)\to 0, 
		\end{equation}
		so the blocks are exponentially separated in scale. Moreover, within each block the distribution is exactly uniform. Thus  $\gamma(\beta)$ has multiscale profile $(q^1,\ldots, q^k)$ with $q^a=\left(\frac{1}{m_a},\ldots,\frac{1}{m_a}\right)\in\Delta_{m_a-1}$ the flat vector for each $a=1,\ldots,k$.
		
		
		As the temperature decreases ($\beta\to\infty$), the Gibbs distribution concentrates on the ground-state sector. More precisely, the limiting Gibbs state is uniform in the last block of dimension $m_k$, and zero everywhere else:
		\begin{equation}
			\label{eq:limit_Gibbs}
			\gamma(\beta)
			\to
			\gamma(\infty)
			=
			\frac1{m_k}\big(
			\underbrace{0,\ldots,0}_{n-m_k},
			\underbrace{1,\ldots,1}_{m_k}
			\big).
		\end{equation}
		The semigroup $\mathcal{S}^{\gamma(\infty)}$ preserving the limiting Gibbs state~\eqref{eq:limit_Gibbs}
		consists of row-stochastic matrices with blocks
		\begin{equation}
			P^{k1}=0, P^{k2}=0,\ldots,  P^{k,{k-1}}=0  
		\end{equation}
		while the lower-right $m_k\times m_k$ block $P^{kk}$ is bistochastic.
	\end{ex}
	\begin{ex} We now give a non-Gibbs example. 
		Let $$\gamma(\beta)=c_{\beta}\left(e^{-\beta}, 2e^{-\beta},\frac{1}{\beta},\frac{1}{\beta},1,2,3\right),$$ where $c_{\beta}=(3e^{-\beta}+\frac{2}{\beta}+6)^{-1}$. This is a vector of type $\mu=(2,2,3)$ as $\beta\to\infty$. Its multiscale profile is 
		$$ 
		q^1=\left(\frac{1}{3},\frac{2}{3}\right),\qquad
		q^2=\left(\frac{1}{2},\frac{1}{2}\right),\qquad
		q^3=\left(\frac{1}{6},\frac{1}{3},\frac{1}{2}\right).
		$$
	\end{ex}
	
	Back to our question: \emph{does $\mathcal{S}^{\gamma(\beta)}$ converge to $\mathcal{S}^{\gamma(\infty)}$, or does a more rigid structure emerge in the limit?}
	
	It turns out that transitions against the blocks of $\gamma(\beta)$ that are separated in scale are suppressed as $\beta\to\infty$. What survives in the limit is therefore not the full  semigroup $\mathcal{S}^{\gamma(\infty)}$, but a subsemigroup of it with a block-upper-triangular structure.
	
	\begin{thm}
		\label{thm:general}
		Let $\left(\gamma(\beta)\right)_{\beta\geq0}$ be a family of probability vectors of type $\mu=(m_1,\ldots,m_k)$ as $\beta\to\infty$, with multiscale profile $(q^1,\ldots,q^k)$. Then the semigroups $\mathcal{S}^{\gamma(\beta)}$
		converge, in the Kuratowski sense, to the set $\mathcal T^{q^1,\ldots,q^k}$ of row-stochastic block-upper-triangular matrices 
		\begin{equation}
			P=
			\begin{pmatrix}
				P^{11} & P^{12} & \cdots & \cdots & P^{1k}\\
				0 & P^{22} & \cdots & & \vdots\\
				\vdots & \vdots & \ddots & & \vdots\\
				\vdots & & & P^{k-1,k-1} & P^{k-1,k}\\
				0 & \cdots & \cdots & 0 & P^{kk}
			\end{pmatrix},
		\end{equation}
		where each block $P^{ab}$ has size $m_a\times m_b$, and the diagonal blocks $P^{aa}$ are $q^a$-subpreserving:
		\begin{equation}
			q^aP^{aa}\leq q^a
		\end{equation}
		for all $a=1,\ldots,k$. (In particular, the last block $P^{kk}$ is $q^k$-preserving.)
		
		As a consequence,  for $n\geq3$, the set-valued map $q\mapsto \mathcal{S}^q$ is not Painlev\'e-Kuratowski continuous at the boundary points of $\Delta_{n-1}$. 
	\end{thm}
	
	\subsection{Zero-temperature limit and BT-stochastic matrices}
	We now specialise Theorem~\ref{thm:general} to Gibbs-preserving semigroups. In this case,  transitions against the energy ordering become exponentially suppressed as $\beta\to\infty$. 
	
	\begin{thm}
		\label{thm:main}
		Let $\mu=(m_1,\ldots,m_k)$ be a composition of $n$, and let
		$$
		E_1\ge\cdots\ge E_n
		$$
		be a vector of type $\mu$. Then the semigroups $\mathcal{S}^{\gamma(\beta)}$
		converge, in the Kuratowski sense, to the set $\mathcal T_\mu$ of row-stochastic block-upper-triangular matrices $P=(P^{ab})_{a,b=1}^k$, 
		where each block $P^{ab}$ has size $m_a\times m_b$, $P^{ab}=0$ if $b<a$, and 
		every diagonal block
		\begin{equation}
			P^{11},P^{22},\ldots,P^{kk}
		\end{equation}
		is column-substochastic.
	\end{thm}
	We call the block-upper-triangular matrices in $\mathcal{T}_{\mu}$, BT-stochastic matrices.
	\begin{rem}
		The last diagonal block $P^{kk}$ is both row-stochastic and column-substochastic. Since its entries are nonnegative, it is automatically column-stochastic, hence bistochastic.
	\end{rem}
	
	\begin{rem}
		In the fully degenerate case $\mu=(n)$,
		$$
		E_1=\cdots=E_n,
		$$
		for all $\beta\geq0$,  $S^{\gamma(\beta)}=\mathcal T_\mu$ is the set of bistochastic matrices. 
		
		Whenever $\mu\neq (n)$, the limit semigroup $\mathcal{T}_{\mu}$ is a proper subsemigroup of $\mathcal{S}^{\gamma(\infty)}$. For example, in the nondegenerate case
		$$
		E_1>\cdots>E_n,
		\qquad
		\mu=(1,\ldots,1),
		$$
		the limiting semigroup $\mathcal T_\mu$ is the set of upper triangular row-stochastic matrices discovered by Narasimhachar and  Gour ~\cite{gour2015resource}. 
	\end{rem}
	
	The difference between $\mathcal{T}_\mu$ and $\mathcal{S}^{\gamma(\infty)}$ can be understood in terms of allowed probability flows between energy levels: $\mathcal{S}^{\gamma(\infty)}$ inhibits only transitions from the ``lowest'' or ``ground'' energy level to the ``excited'' ones, but still allows transitions from lower to upper energy levels among the excited ones, while this is not possible within $\mathcal{T}_\mu$, which is reminiscent of the level hierarchy in the low-temperature Gibbs weights.
	Informally, if a row-stochastic matrix $P$ is in the $\beta\to\infty$ limit of the family $\{\mathcal{S}^{\gamma(\beta)}\}_{\beta>0}$, then there exist matrices  $P^{\beta}\in \mathcal{S}^{\gamma(\beta)}$, such that $P^{\beta}\to P$ as $\beta\to\infty$. This implies that 
	\begin{equation}
		\gamma(\beta)(P-I)=\gamma(\beta)(P-P^{\beta})+\underbrace{\gamma(\beta)(P^{\beta}-I)}_{=0}= \gamma(\beta) (P-P^{\beta})\to 0.
	\end{equation}
	But by continuity, $\gamma(\beta)(P-I)\to \gamma(\infty)(P-I)$. We conclude that $\gamma(\infty)P=\gamma(\infty)$, namely $P\in \mathcal{S}^{\gamma(\infty)}$. 
	
	However, not all matrices $P\in \mathcal{S}^{\gamma(\infty)}$ can be approximated by $\gamma(\beta)$-preserving stochastic matrices for large $\beta$. For instance, consider 
	\begin{equation}
		P=
		\mqty(1 & 0 &0\\
		1 &0 &0 \\
		0 & 0 &1)\in \mathcal{S}^{\gamma(\infty)}.
	\end{equation}
	Assume that there exists a family $\{P^{\beta}\}_{\beta}$ with $P^{\beta}\in \mathcal{S}^{\gamma(\beta)}$ such that $\|P^{\beta}-P\|\to0$, as $\beta\to\infty$. For all $\beta>0$, since $\gamma(\beta)P^{\beta}=\gamma(\beta)$, we have 
	\begin{equation}
		\gamma_1(\beta)P^{\beta}_{1,1}
		+\gamma_2(\beta)P^{\beta}_{2,1}
		+ \gamma_3(\beta)P^{\beta}_{3,1}=\gamma_1(\beta).
	\end{equation}
	Hence, $\gamma_2(\beta)P^{\beta}_{2,1}\leq \gamma_1(\beta)$, and we get to the contradiction
	$$
	P^{\beta}_{2,1}\leq \frac{\gamma_1(\beta)}{\gamma_2(\beta)}\to 0\neq 1=P_{2,1}.
	$$
	
	Indeed, Theorem~\ref{thm:main} claims that the zero-temperature limit of $\mathcal{S}^{\gamma(\beta)}$ is the set of row-stochastic upper triangular matrices $\mathcal{T}_{(1,1,1)}$, a proper subsemigroup of $ \mathcal{S}^{\gamma(\infty)}$.
	
	\subsection{Extreme points of BT-stochastic matrices}
	
	Recall that the set $S^q$ of $q$-preserving stochastic matrices is compact and convex. By a classical theorem of Minkowski~\cite{brondsted2012introduction}, every compact convex set is the convex hull of its set of extreme points. For general $q$, the characterisation of the extreme points of $S^q$ is non-trivial. It is known~\cite[Remark 4.5]{joe1990majorization} that, for $q\in\Delta_{n-1}$,
	the number of extreme points of $S^q$ is at least $n!$ and at most $\binom{n^2}{2n-1}$. 
	If $q=\e$, then  $S^q$ is the celebrated Birkhoff polytope whose extreme points are the $n!$ permutation matrices~\cite{birkhoff1946tres}. 
	
	
	
	It is easy to see that $\mathcal{T}_{\mu}$ is also a convex set whose extreme points can be characterised and enumerated. 
	
	\begin{prop}
		\label{prop:extreme_pts}
		Let $\mu=(m_1,\ldots, m_k)$ be a composition of $n$.  Denote by $\operatorname{ext}(\mathcal{T}_{\mu})$ the set of extreme points of $\mathcal{T}_{\mu}$. Then, 
		\begin{enumerate}
			\item $\operatorname{ext}(\mathcal{T}_{\mu})$ is the set of  block-upper-triangular $0$--$1$ matrices satisfying:
			\begin{enumerate}
				\item each row contains exactly one $1$;
				\item for every $a=1,\ldots,k$, each column of the diagonal block $P_{aa}$ contains at most one $1$.
			\end{enumerate}
			\item The number of extreme points $\operatorname{ext}(\mathcal{T}_{\mu})$ is
			\begin{equation}
				\label{eq:Laguerre_repr}
				\#\operatorname{ext}(\mathcal{T}_{\mu})=\prod_{l=1}^k\widehat{L}_{m_l}\left(m_1+\cdots+m_{l-1}\right),
			\end{equation}
			where $\widehat{L}_n(x):=n!L_n(-x)$, with
			\begin{equation}
				L_n(x):=[t^n]\frac{1}{1-x}e^{-\frac{tx}{1-x}}=\sum_{k=0}^{n}\binom{n}{k}\frac{(-1)^k}{k!}x^k
			\end{equation} 
			the $n$-th Laguerre polynomial, 
			%$m_{<l}:=m_1+\cdots+m_{l-1}$, 
			and we use the convention that the empty sum is $0$.
		\end{enumerate} 
		
	\end{prop}
	
	\subsection{BT-majorisation}
	The semigroup $\mathcal{T}_{\mu}$ defines a semigroup majorisation on probability vectors~\cite{Cunden2025SemigroupMajorisation}. For $x,y\in \R^n$ we say that $x$ is $\mathcal{T}_{\mu}$-majorised by $y$, and write $x\prec^{\mathcal{T}_{\mu}}y$ if $$
	x=yP,\, \text{ for some $P\in\mathcal{T}_{\mu}$.}
	$$
	The relation $\prec^{\mathcal{T}_{\mu}}$ is a preorder on $\R^n$. We call this relation BT-majorisation.
	
	We index vectors $x\in\R^n$ according to the block structure  $\mu=(m_1,\ldots,m_k)$, writing $x=(x^1,\ldots,x^k)$ where $x^a=(x^a_1,\ldots,x^a_{m_a})$ for all $a=1,\ldots,k$.
	
	For each $a=1,\ldots,k$,  let $(x^a)^{\downarrow}$ be the nonincreasing rearrangement of the entries $x^a$. The $\mu$-nonincreasing rearrangement of $x$ is then defined by 
	\begin{equation}\label{eq:NotationMuOrdering}
		x^{\mu,\downarrow}=\left((x^1)^{\downarrow},(x^2)^{\downarrow},\ldots, (x^k)^{\downarrow}\right).
	\end{equation}
	In words, $x^{\mu,\downarrow}$ is the vector obtained by sorting the entries of $x$ 
	within the blocks, while preserving the block order.
	
	We also define the \textbf{lexicographic} $<
	^{\mathrm{lex}}$ and \textbf{reverse lexicographic} $<
	^{\mathrm{revlex}}$ orders. For $a=(a_1,\ldots,a_k)$ and $b=(b_1,\ldots,b_k)$ real $n$-tuples, we write  $a<
	^{\mathrm{lex}}b$ if, at the first differing coordinate $i$, $a_i<b_i$; we write instead $a<
	^{\mathrm{revlex}}b$ if, at the last differing coordinate $i$, $a_i<b_i$.
	
	
	\begin{thm}
		\label{prop:monotones}
		Let $\mu$ be a composition of $n$.
		For $x,y\in\mathbb \R_+^n$, the following are equivalent:
		\begin{enumerate}
			\item $x\prec^{\mathcal T_\mu}y$;
			\item for every $\ell=1,\ldots,n$,
			$$
			\sum_{i=1}^{\ell}x_i^{\mu,\downarrow}
			\le
			\sum_{i=1}^{\ell}y_i^{\mu,\downarrow},
			$$
			with equality for $\ell=n$.
		\end{enumerate}
	\end{thm}
	
	BT-majorisation interpolates between standard majorisation (when $\mu=(n)$) and upper triangular (UT) or unordered majorisation (when $\mu=(1,\ldots,1)$).
	\subsection{The fate of monotones}
	It is known~\cite{vom2019d} that
	\begin{equation}
		\label{eq:thermo}
		x\prec^{\gamma(\beta)} y\quad\iff\quad \sum_j \gamma(\beta)_j\, g\!\left(\frac{x_j}{\gamma(\beta)_j}\right)
		\le
		\sum_j \gamma(\beta)_j\, g\!\left(\frac{y_j}{\gamma(\beta)_j}\right)
	\end{equation}
	for every continuous convex function $g$ defined on a domain containing the points
	$
	x_j/\gamma(\beta)_j
	$
	and
	$
	y_j/\gamma(\beta)_j
	$. Assume that $\gamma(\beta)$ is a Gibbs vector parametrised by the inverse temperature $\beta$.
	
	Choosing
	$$
	g(t)=\frac{1}{\beta}t\log t,
	\qquad t\ge0,
	$$
	in~\eqref{eq:thermo},
	one finds, 
	\begin{equation}
		x\prec^{\gamma(\beta)} y
		\;\Longrightarrow\;
		F_\beta(x)\le F_\beta(y),
	\end{equation}
	where
	$$
	F_\beta=U-\frac1\beta S
	$$
	is the free energy, 
	\begin{equation}
		U(x)=\sum_j x_j E_j,
		\qquad
		S(x)=-\sum_j x_j\log x_j
	\end{equation}
	with $U$ and $S$ denoting the mean energy and entropy, respectively.
	
	Choosing instead
	\begin{equation}
		g(t)
		=
		\frac{1}{\alpha-1}t^\alpha,
		\qquad
		\alpha>0,\,\, \alpha\neq 1,
	\end{equation}
	which is strictly convex on $t>0$, one obtains
	\begin{equation}
		x\prec^{\gamma(\beta)} y
		\;\Longrightarrow\;
		P_\alpha(x)\le P_\alpha(y),
	\end{equation}
	where
	\begin{equation}
		P_\alpha(x)
		=
		\frac{1}{\alpha-1}
		\frac{1}{Z(\beta)^{\alpha-1}}
		\sum_j x_j^\alpha e^{-\beta(1-\alpha)E_j},
	\end{equation}
	is a  linearised relative $\alpha$-entropy.
	
	In the language of physics, if $x$ is thermomajorised by $y$ at temperature $1/\beta$, then the free energy and the $\alpha$-entropies of $x$ do not exceed those of $y$. The following proposition shows that, in the zero-temperature limit, these entropic inequalities collapse to much simpler conditions.
	\begin{prop}
		\label{prop:fate_monotones}
		Let $x,y\in\Delta_{n-1}$ and suppose that $x\prec^{\gamma(\beta)} y$  for all sufficiently large $\beta$.  Then
		either
		$$
		U(x)<U(y),
		$$
		or
		$$
		U(x)=U(y)
		\quad\text{and}\quad
		S(x)\ge S(y).
		$$
		Moreover, 
		\begin{equation}
			\label{eq:revlex-condition}
			\left(
			\|x^1\|_1,
			\dots,
			\|x^k\|_1
			\right)
			<
			^{\mathrm{revlex}}
			\left(
			\|y^1\|_1,
			\dots,
			\|y^k\|_1
			\right),
		\end{equation}
		and
		\begin{equation}
			\label{eq:lex-condition}
			\left(
			\|y^1\|_1,
			\dots,
			\|y^k\|_1
			\right)
			<
			^{\mathrm{lex}}
			\left(
			\|x^1\|_1,
			\dots,
			\|x^k\|_1
			\right),
		\end{equation} 
		where $\|x^a\|_1=\sum_{j=1}^{m_a}x^a_j$ and $\|y^a\|_1=\sum_{j=1}^{m_a}y^a_j$, for all $a=1,\ldots,k$.
	\end{prop}
	Thus, in the low-temperature limit, the mean energy is the primary monotone in the free energy, while entropy provides the subleading order within a fixed energy sector.
	Indeed, if $U(x)=U(y)$, then
	\begin{equation}
		F_\beta(y)-F_\beta(x)
		=
		\frac1\beta\bigl(S(x)-S(y)\bigr),
	\end{equation}
	so transitions preserving the mean energy can only increase entropy. Restricted to a fixed-energy subspace, the dynamics reduces to bistochastic mixing, which drives states toward the uniform distribution on that subspace. Decreasing entropy therefore requires population transfer to lower energy sectors.
	
	The $\alpha$-entropic inequalities all collapse to a reverse lexicographic condition on the block sums when $1-\alpha>0$, and a lexicographic condition when $1-\alpha<0$.
	
	\section{Discussion of Gibbs-preserving matrices in low dimension} \label{sec:examples}
	
	\subsection{BT-matrices in dimensions $n=3$ and $n=4$}
	
	%\begin{ex}
	For $n=3$, there are four possible types of degeneracies of 
	$E_1 \ge E_2 \ge E_3 $:
	\begin{table}[h]
		\centering
		\begin{tabular}{r|c|c|c}
			$\mu$ &  $\#$ blocks $k$ & energy structure & $\#\operatorname{ext}(\mathcal{T}_{\mu})$ \\
			\hline
			$(1,1,1)$ & 3 & $E_1>E_2>E_3$ & 6\\
			$(1,2)$ & 2 & $E_1>E_2=E_3$& 6\\
			$(2,1)$ & 2 & $E_1=E_2>E_3$& 7\\
			$(3)$ & 1 & $E_1=E_2=E_3$ & 6 
		\end{tabular}
	\end{table}
	
	
	We now describe the limiting semigroups $\mathcal T_\mu$.
	The pattern is always the same: blocks above the energy ordering survive, blocks below vanish, and diagonal blocks carry the column-substochastic constraints. The last diagonal block is bistochastic.
	\begin{align}
		\mathcal T_{(1,1,1)}&=
		\left\{
		\begin{pmatrix}
			* & * & *\\
			0 & * & * \\
			0 & 0 & * \\
		\end{pmatrix}
		\text{ row-stochastic}
		\right\},\\
		\mathcal T_{(1,2)}&=
		\left\{
		\begin{pmatrix}
			* & * & * \\
			0 & \square & \square\\
			0 & \square & \square
		\end{pmatrix}
		\text{ row-stochastic}
		\colon
		\begin{pmatrix}
			\square  & \square \\
			\square  & \square 
		\end{pmatrix}
		\text{ bistochastic}
		\right\},\nonumber\\
		\mathcal T_{(2,1)}&=
		\left\{
		\begin{pmatrix}
			\triangle & \triangle  & *\\
			\triangle & \triangle & *\\
			0 & 0 & 1
		\end{pmatrix}
		\text{ row-stochastic}
		\colon
		\begin{pmatrix}
			\triangle & \triangle \\
			\triangle & \triangle \\
		\end{pmatrix}
		\text{ column-substochastic}
		\right\},\nonumber\\
		\mathcal T_{(3)}&=\left\{
		\begin{pmatrix}
			\square & \square  & \square\\
			\square & \square & \square\\
			\square & \square & \square
		\end{pmatrix} \text{ 
			bistochastic}\right\}.\nonumber
	\end{align}
	
	For $n=4$, there are eight possible types of degeneracies of 
	$E_1 \ge E_2 \ge E_3 \ge E_4$:
	
	\begin{table}[h]
		\centering
		\begin{tabular}{r|c|c|c}
			$\mu$ & Number of blocks $k$ & energy structure &  $\#\operatorname{ext}(\mathcal{T}_{\mu})$ \\
			\hline
			$(1,1,1,1)$ & 4 & $E_1>E_2>E_3>E_4$ & 24\\
			$(1,1,2)$ & 3 & $E_1>E_2>E_3=E_4$ &24\\
			$(1,2,1)$ & 3 & $E_1>E_2=E_3>E_4$ &28\\
			$(1,3)$ & 2 & $E_1>E_2=E_3=E_4$ & 24\\
			$(2,1,1)$ & 3 & $E_1=E_2>E_3>E_4$ & 28\\
			$(2,2)$ & 2 & $E_1=E_2>E_3=E_4$ & 28\\
			$(3,1)$ & 2 & $E_1=E_2=E_3>E_4$ & 34\\
			$(4)$ & 1 & $E_1=E_2=E_3=E_4$ & 24
		\end{tabular}
	\end{table}
	
	The limiting semigroups $\mathcal T_\mu$ are:
	\begin{align}
		% \begin{aligned}
			\mathcal T_{(1,1,1,1)}&=
			\left\{
			\begin{pmatrix}
				* & * & * & *\\
				0 & * & * & *\\
				0 & 0 & * & *\\
				0 & 0 & 0 & *
			\end{pmatrix}
			\text{ row-stochastic}
			\right\},\\
			\mathcal T_{(1,1,2)}&=
			\left\{
			\begin{pmatrix}
				* & * & * & *\\
				0 & * & * & *\\
				0 & 0 & \square & \square\\
				0 & 0 & \square & \square
			\end{pmatrix}
			\text{ row-stochastic}
			\colon
			\begin{pmatrix}
				\square  & \square \\
				\square  & \square 
			\end{pmatrix}
			\text{  bistochastic}
			\right\},\nonumber\\
			\mathcal T_{(1,2,1)}&=
			\left\{
			\begin{pmatrix}
				* & * & * & *\\
				0 & \triangle & \triangle & *\\
				0 & \triangle & \triangle & *\\
				0 & 0 & 0 & 1
			\end{pmatrix}
			\text{ row-stochastic}
			\colon
			\begin{pmatrix}
				\triangle & \triangle\\
				\triangle & \triangle
			\end{pmatrix}
			\text{ column-substochastic}
			\right\},\nonumber\\
			\mathcal T_{(1,3)}&=
			\left\{
			\begin{pmatrix}
				* & * & * & *\\
				0 & \square &  \square &  \square\\
				0 & \square & \square &  \square\\
				0 &  \square &  \square &  \square
			\end{pmatrix}
			\text{ row-stochastic}
			\colon
			\begin{pmatrix}
				\square &  \square &  \square\\
				\square &  \square &  \square\\
				\square &  \square &  \square
			\end{pmatrix}
			\text{  bistochastic}
			\right\},\nonumber\\
			\mathcal T_{(2,1,1)}&=
			\left\{
			\begin{pmatrix}
				\triangle & \triangle & * & *\\
				\triangle & \triangle & * & *\\
				0 & 0 & * & *\\
				0 & 0 & 0 & 1
			\end{pmatrix}
			\text{ row-stochastic}
			\colon
			\begin{pmatrix}
				\triangle & \triangle\\
				\triangle & \triangle
			\end{pmatrix}
			\text{ column-substochastic}
			\right\},\nonumber\\
			\mathcal T_{(2,2)}&=
			\left\{
			\begin{pmatrix}
				\triangle & \triangle  & * & *\\
				\triangle  & \triangle  & * & *\\
				0 & 0 & \square & \square\\
				0 & 0 & \square & \square
			\end{pmatrix}
			\text{ row-stochastic}
			\colon
			\begin{aligned}
				&
				\begin{pmatrix}
					\triangle & \triangle\\
					\triangle & \triangle
				\end{pmatrix} \text{ column-substochastic},\\
				&
				\begin{pmatrix}
					\square & \square \\
					\square  & \square 
				\end{pmatrix} \text{ bistochastic}
			\end{aligned}
			\right\},\nonumber\\
			\mathcal T_{(3,1)}&=
			\left\{
			\begin{pmatrix}
				\triangle & \triangle & \triangle & *\\
				\triangle & \triangle & \triangle & *\\
				\triangle & \triangle & \triangle & *\\
				0 & 0 & 0 & 1
			\end{pmatrix}
			\text{ row-stochastic}
			\colon
			\begin{pmatrix}
				\triangle & \triangle & \triangle\\
				\triangle & \triangle & \triangle\\
				\triangle & \triangle & \triangle
			\end{pmatrix}
			\text{ column-substochastic}
			\right\},\nonumber\\
			\mathcal T_{(4)}&=\left\{
			\begin{pmatrix}
				\square & \square & \square & \square\\
				\square & \square & \square & \square\\
				\square & \square & \square & \square\\
				\square & \square & \square & \square
			\end{pmatrix} \text{ bistochastic}\right\}.\nonumber
		\end{align}
		
		
		\subsection{Finite-temperature corrections} 
		
		The zero-temperature limit captures the leading-order behaviour of Gibbs-preserving transformations at low temperatures in the same way that permutations capture leading-order behaviour at high temperatures. A natural next step is to understand finite-temperature corrections: can one describe asymptotic expansions of the semigroup and of the induced preorder in powers of $e^{-\beta}$? Such corrections would reveal how BT-majorisation deforms into thermomajorisation. 
		
		It is well known that at infinite temperature the extreme operations are precisely permutations, i.e. the vertices of the Birkhoff polytope. This statement is independent of the energy type of the system. At the opposite end, in the nondegenerate zero-temperature limit, the extremal operations are given by upper triangular stochastic matrices. The BT-stochastic matrices introduced above close this gap at the level of the limiting theory: for a Gibbs family of type $\mu=(m_1,\ldots,m_k)$, they give the limiting semigroup $T_\mu$, and hence an exact description of the extreme operations which survive as $\beta\to\infty$. Away from these limiting regimes, however, the structure of extreme operations is  harder to characterise.
		
		A useful perspective on the emergence of finite-temperature extremal operations is provided by the linear-programming construction of Ref.~\cite{mazurek2018decomposability}. Alternatively, one can use construction of Jurkat and Ryser~\cite{Jurkat1967}. In the zero- and infinite-temperature regimes, the vertices of the corresponding polytopes are known explicitly. One may then follow individual extremal branches into the finite-temperature region by perturbing a selected extremal operation while preserving both stochasticity and the stationary state. Extremality is maintained by saturating the maximal number of admissible positivity constraints, which amounts to selecting subsets of matrix entries that remain fixed at zero. Repeating this procedure over all admissible choices generates the full graph of extremal branches emanating from the limiting vertices. 
		
		Not all such branches persist throughout the entire temperature range. As the temperature changes, positivity may fail at isolated critical inverse temperatures, leading to the coalescence or branching of extremal operations. These critical points occur whenever
		%
		% One obstruction is the possible appearance of critical inverse temperatures, at which 
		two partial sums of the invariant state coincide. In the Gibbs case, for two subsets $A,B\subset\{1,\ldots,n\}$, after removing common indices, such critical points are solutions of 
		
		\begin{equation}
			\sum_{a\in A}\gamma_a(\beta) = \sum_{b\in B}\gamma_b(\beta) \quad\Longleftrightarrow\quad
			\sum_{a\in A}e^{-\beta E_a} = \sum_{b\in B}e^{-\beta E_b}, 
		\end{equation}
		since the Gibbs normalisation cancels. 
		
		The argument which bounds the possible critical temperatures, however, does not rely on the exponential form itself. What matters is the multiscale structure of the invariant distribution. We therefore formulate the estimate for a general family of type $\mu=(m_1,\ldots,m_k)$ in the sense of Section~\ref{sec:multiscale}. To this end we define the block masses 
		
		\begin{equation}
			w_a(\beta):=\|\gamma^a(\beta)\|_1 . 
		\end{equation}
		
		Let $I_a$ denote the set of indices in the $a$-th block. For disjoint $A,B\subset\{1,\ldots,n\}$, define 
		
		\begin{equation}
			d_a^{A,B}(\beta) := \frac{1}{w_a(\beta)}\qty(\sum_{i\in A\cap I_a} \gamma_i(\beta) - \sum_{i\in B\cap I_a} \gamma_i(\beta)) . 
		\end{equation}
		Then the criticality condition becomes 
		
		\begin{equation}
			\sum_{a=1}^k d_a^{A,B}(\beta)\,w_a(\beta)=0 . 
		\end{equation}
		
		Assume, in addition, that for every $a<r$ the relative scale functions 
		
		\begin{equation}
			\rho_{ar}(\beta) := \frac{w_a(\beta)}{w_r(\beta)} 
		\end{equation}
		are continuous, non-increasing in $\beta$, and satisfy $\rho_{ar}(\beta)\to0$ as $\beta\to\infty$. Suppose that the criticality equation is nontrivial, and let $r$ be the largest block index for which $d_r^{A,B}\neq0$. Changing the roles of $A$ and $B$ if necessary, assume $d_r^{A,B}>0$. Dividing by $w_r(\beta)$ gives 
		
		\begin{equation}
			d_r^{A,B}(\beta) = -\sum_{a<r} d_a^{A,B}(\beta)\, \rho_{ar}(\beta). 
		\end{equation}
		Hence, for all sufficiently large $\beta$, 
		
		\begin{equation}\label{eq:gen_max_crit}
			0<c_r \leq \sum_{a<r} C_a\,\rho_{ar}(\beta) \equiv R_r(\beta), 
		\end{equation}
		where $c_r>0$ is a lower bound for $d_r^{A,B}(\beta)$ and the constants $C_a$ bound the possible absolute contribution from the higher-energy blocks. Since $R_r(\beta)\to0$, no nontrivial critical inverse temperature associated with the block $r$ can occur beyond the threshold determined by 
		
		\begin{equation}
			R_r(\beta_r)=c_r, 
		\end{equation}
		whenever such a solution exists. If the functions $R_r$ are strictly decreasing, this threshold is unique. We now specialise this general multiscale estimate back to the Gibbs family. Let $E_1>\cdots>E_k$ be the distinct energy values, ordered as in the rest of the paper, with multiplicities $m_1,\ldots,m_k$. In the Gibbs case each block profile is flat, which yields
		
		\begin{equation}
			w_a(\beta) = \frac{m_a e^{-\beta E_a}}{Z(\beta)}\qq{and} \rho_{ar}(\beta) = \frac{m_a}{m_r} e^{-\beta(E_a-E_r)}.
		\end{equation}
		Equivalently, multiplying out the flat block factors, the criticality condition can be written as 
		
		\begin{equation}
			\sum_{a=1}^k c_a e^{-\beta E_a}=0, 
		\end{equation}
		where $
		c_a := \abs{\{a\in A:E_a=E_a\}} - \abs{\{b\in B:E_b=E_a\}}$.
		If all $c_a$ vanish, the equality is an identity caused solely by degeneracies and does not define an isolated critical temperature. Otherwise, Eq. \eqref{eq:gen_max_crit} turns into
		
		\begin{equation}
			1 \leq c_r \leq \sum_{a<r} \abs{c_a}e^{-\beta(E_a-E_r)} \leq \sum_{a<r} m_a e^{-\beta(E_a-E_r)} \equiv R_r(\beta). 
		\end{equation}
		Here $E_a>E_r$ for $a<r$, hence $R_r(\beta)$ is continuous and strictly decreasing whenever there is at least one higher energy level, with $\lim_{\beta\to\infty}R_r(\beta)=0$. Therefore, if $R_r(0)>1$, there is a unique positive solution $\beta_r$ of 
		
		\begin{equation*}
			R_r(\beta_r)=1 . 
		\end{equation*}
		Every nontrivial Gibbs critical inverse temperature $\beta_*$ must therefore satisfy 
		
		\begin{equation}
			\label{eq:minimal_crit_temp} \beta_* \leq \max_r \beta_r =: \beta_{\rm low}. 
		\end{equation}
		Conversely, each such threshold is realised by taking $A$ to contain a single state at energy $E_r$, and $B$ to contain all states with energies strictly larger than $E_r$. Thus, in the Gibbs case, the abstract multiscale obstruction has a simple exponential form: a selected population in the block $E_r$ is balanced against the cumulative Gibbs weight of higher-energy blocks. At such a critical temperature, additional extremal operations of the type described in \cite{mazurek2018decomposability} appear. In the zero-temperature limit these coincidences disappear, and the only extremal structure which survives is the BT-stochastic one encoded by $T_\mu$.
		
		% \blue{
			The sequence of critical temperatures also suggests a constructive procedure for determining the finite-temperature extremal structure. Let 
			$$ \beta_{\rm low}=: \beta_0 > \beta_1 > \cdots > \beta_N := \beta_{\rm high} $$ 
			denote the ordered list of all critical inverse temperatures. 
			
			Starting from the low-temperature regime, one first determines the extremal BT-stochastic operations associated with the semigroup $\mathcal{T}_\mu$. These serve as root vertices for the subsequent construction. One then proceeds through the temperature intervals $(\beta_{j+1},\beta_j)$ in succession: 
			
			\begin{enumerate} 
				\item At a critical point $\beta_j$, identify all candidate extremal operations whose positivity constraints become saturated and which cease to be positive for $\beta>\beta_j$. 
				
				\item Treat these operations as root vertices at $\beta=\beta_j$ and apply the linear-programming construction of Ref.~\cite{mazurek2018decomposability} to generate all admissible extremal branches on the interval $(\beta_{j+1},\beta_j)$. 
				
				\item Within the interval, follow each branch continuously while maintaining stochasticity and stationarity with respect to the Gibbs state. 
				
				\item Determine which branches terminate at the next critical temperature $\beta_{j+1}$ and which remain positive throughout the interval. 
				
				\item Use all operations becoming critical at $\beta_{j+1}$ as the new collection of root vertices and repeat the procedure. 
			\end{enumerate} 
			Iterating this construction from $\beta_0$ down to $\beta_N$ generates, in principle, the complete graph of extremal Gibbs-preserving operations across the entire temperature range. The resulting structure may then be verified by performing the same construction in the opposite direction, starting from the permutation vertices governing the infinite-temperature limit and propagating the extremal branches towards lower temperatures.
			
			
			In particular, for $n=3$ there is only one  critical temperature $\beta_*$ defined by $\exp[-\beta(E_2 - E_3)] + \exp[-\beta(E_1 - E_3)] = 1$, thus defining two regimes $\beta \leq \beta_*$ and $\beta \geq \beta_*$. 
			Complete code used to generate extreme operations is available at \cite{github}; the code is restricted to roots at very high and very low temperature regimes. Resulting branching of extreme operations is summarised in graphs shown in Figs. \ref{fig:flowchart_21_12} and \ref{fig:flowchart_111}.
			%
			Throughout the following graphs we will use convention
			\begin{equation}
				\begin{aligned}
					P_1 & = \mqty(
					1 & 0 & 0 \\
					0 & 1 & 0 \\
					0 & 0 & 1), & 
					P_2 & = \mqty(
					1 & 0 & 0 \\
					0 & 0 & 1 \\
					0 & 1 & 0), &
					P_3 & = \mqty(
					0 & 1 & 0 \\
					1 & 0 & 0 \\
					0 & 0 & 1), \\
					P_4 & = \mqty(
					0 & 0 & 1 \\
					1 & 0 & 0 \\
					0 & 1 & 0
					), & 
					P_5 & = \mqty(
					0 & 1 & 0 \\
					0 & 0 & 1 \\
					1 & 0 & 0
					), &
					P_6 & = \mqty(
					0 & 0 & 1 \\
					0 & 1 & 0 \\
					1 & 0 & 0
					).
				\end{aligned}
			\end{equation}
			
			We begin with the case where degeneracy happens among the excited levels, corresponding to the $(2,1)$ structure $E_2 = E_1$, thus prompting the critical temperature to be given by
			\begin{equation}
				\beta_* = \frac{\log(2)}{E_1 - E_3} 
			\end{equation}
			where $E_3$ is the ground state energy and $E_1$ is the excited level energy. This provides us with the extreme operation flowchart as presented in the left panel of Fig. \ref{fig:flowchart_21_12}.
			
			\begin{figure}[h]
				\centering
				\includegraphics[width=.475\linewidth]{21_structure.pdf}
				\hspace{.02\linewidth}
				\includegraphics[width=.475\linewidth]{12_structure.pdf}
				\caption{Flowchart for extreme points of $\mathcal{S}^{\gamma(\beta)}$ for $n=3$ under level degeneracy with $\mu = (2,1)$ \textit{[left panel]} and $\mu = (1,2)$ \textit{[right panel]}, corresponding to doubly degenerate excited and ground levels, respectively. Definitions of $A^{(2,1)}_i$ and $A^{(1,2)}_i$ operations are given in equations \eqref{eq:21_ext_ops} and \eqref{eq:12_ext_ops}, respectively.}
				\label{fig:flowchart_21_12}
			\end{figure}
			We note that $P_1 = A_1^{(2,1)} = T_1^{(2,1)}$ and $P_2 = A_2^{(2,1)} = T_2^{(2,1)}$ as both represent valid permutations restricted to the excited level subset. Below we list all extreme operations
			\begin{equation}\label{eq:21_ext_ops}
				{\scriptsize
					\begin{aligned}
						A^{(2,1)}_3 & = \left(
						\begin{array}{ccc}
							1-\frac{\gamma_1}{\gamma_3} & \frac{\gamma_1}{\gamma_3} & 0 \\
							1 & 0 & 0 \\
							0 & 0 & 1 \\
						\end{array}
						\right) &
						A^{(2,1)}_4 & = \left(
						\begin{array}{ccc}
							0 & \frac{\gamma_1}{\gamma_3} & 1-\frac{\gamma_1}{\gamma_3} \\
							1 & 0 & 0 \\
							\frac{\gamma_3}{\gamma_1}-1 & 0 & 2-\frac{\gamma_3}{\gamma_1} \\
						\end{array}
						\right) & 
						A^{(2,1)}_5 & = \left(
						\begin{array}{ccc}
							1-\frac{\gamma _1}{\gamma _3} & 0 & \frac{\gamma _1}{\gamma _3} \\
							1 & 0 & 0 \\
							0 & 1 & 0 \\
						\end{array}
						\right) \\
						A^{(2,1)}_6 & = \left(
						\begin{array}{ccc}
							0 & 1-\frac{\gamma_1}{\gamma_3} & \frac{\gamma_1}{\gamma_3} \\
							1 & 0 & 0 \\
							\frac{\gamma_3}{\gamma_1}-1 & 2-\frac{\gamma_3}{\gamma_1} & 0 \\
						\end{array}
						\right) &
						A^{(2,1)}_7 & = \left(
						\begin{array}{ccc}
							0 & \frac{\gamma_1}{\gamma_3} & 1-\frac{\gamma_1}{\gamma_3} \\
							\frac{\gamma_3}{\gamma_1}-1 & 0 & 2-\frac{\gamma_3}{\gamma_1} \\
							1 & 0 & 0 \\
						\end{array}
						\right) &
						A^{(2,1)}_8 & = \left(
						\begin{array}{ccc}
							1-\frac{\gamma_1}{\gamma_3} & \frac{\gamma_1}{\gamma_3} & 0 \\
							0 & 0 & 1 \\
							1 & 0 & 0 \\
						\end{array}
						\right) \\ 
						A^{(2,1)}_9 & = \left(
						\begin{array}{ccc}
							0 & 1-\frac{\gamma_1}{\gamma_3} & \frac{\gamma_1}{\gamma_3} \\
							\frac{\gamma_3}{\gamma_1}-1 & 2-\frac{\gamma_3}{\gamma_1} & 0 \\
							1 & 0 & 0 \\
						\end{array}
						\right) &
						A^{(2,1)}_{10} &  = \left(
						\begin{array}{ccc}
							1-\frac{\gamma_1}{\gamma_3} & 0 & \frac{\gamma_1}{\gamma_3} \\
							0 & 1 & 0 \\
							1 & 0 & 0 \\
						\end{array}
						\right) & 
						A^{(2,1)}_{11} & = \left(
						\begin{array}{ccc}
							1-\frac{2 \gamma_1}{\gamma_3} & \frac{\gamma_1}{\gamma_3} & \frac{\gamma_1}{\gamma_3} \\
							1 & 0 & 0 \\
							1 & 0 & 0 \\
						\end{array}
						\right).
				\end{aligned}}
			\end{equation}
			The subset $\qty{A_1^{(2,1)},\hdots,A_{10}^{(2,1)}}$ provides the high-temperature limit of the operations existing for $\beta \leq \beta_*$, while restriction to $\qty{A_1^{(2,1)}, A_2^{(2,1)}, A_3^{(2,1)}, A_5^{(2,1)}, A_8^{(2,1)}, A_{10}^{(2,1)}, A_{11}^{(2,1)}}$ is a valid set of extreme points for $\beta > \beta_*$, with each of them giving rise to one of the extreme operations in $\beta\rightarrow\infty$ limit, which can be read off immediately by setting $\gamma_e = 0$ and $\gamma_g = 1$. At the same time, let us note that at $\beta = \beta_*$
			\begin{equation}
				A_4^{(2,1)} = A_6^{(2,1)} = A_7^{(2,1)} = A_9^{(2,1)} = A_{11}^{(2,1)} = C^{(2,1)} = \mqty(
				0 & \frac{1}{2} & \frac{1}{2} \\ 
				1 & 0 & 0 \\ 
				1 & 0 & 0)
			\end{equation}
			
			Next, we proceed to the case of ground state degeneracy $E_2 = E_3$ or $\mu = (1,2)$. This  structure leads to vanishing of the critical temperature, resulting only in splittings at intermediate temperatures, shown in the flow graph in the right panel of Fig.~\ref{fig:flowchart_21_12}.
			
			In this case persistent permutations correspond to $P_1 = A_1^{(1,2)} = T_1^{(1,2)}$ and $P_3 = A_4^{(1,2)} = T_3^{(1,2)}$, with the remaining operations given by
			
			\begin{equation}\label{eq:12_ext_ops}
				{\scriptsize
					\begin{aligned}
						A^{(1,2)}_2 & = \left(
						\begin{array}{ccc}
							1 & 0 & 0 \\
							0 & 1-\frac{\gamma_1}{\gamma_3} & \frac{\gamma_1}{\gamma_3} \\
							0 & 1 & 0 \\
						\end{array}
						\right) & 
						A^{(1,2)}_3 & = \left(
						\begin{array}{ccc}
							\frac{\gamma_1}{\gamma_3} & 1-\frac{\gamma_1}{\gamma_3} & 0 \\
							1-\frac{\gamma_1}{\gamma_3} & 0 & \frac{\gamma_1}{\gamma_3} \\
							0 & 1 & 0  
						\end{array}
						\right) \\
						A^{(1,2)}_5 & = \left(
						\begin{array}{ccc}
							1-\frac{\gamma_1}{\gamma_3} & 0 & \frac{\gamma_1}{\gamma_3} \\
							\frac{\gamma_1}{\gamma_3} & 1-\frac{\gamma_1}{\gamma_3} & 0 \\
							0 & 1 & 0 \\
						\end{array}
						\right) &
						A^{(1,2)}_6 & = \left(
						\begin{array}{ccc}
							0 & 1-\frac{\gamma_1}{\gamma_3} & \frac{\gamma_1}{\gamma_3} \\
							1 & 0 & 0 \\
							0 & 1 & 0 \\
						\end{array}
						\right) & 
						A^{(1,2)}_7 & = \left(
						\begin{array}{ccc}
							1-\frac{\gamma_1}{\gamma_3} & \frac{\gamma_1}{\gamma_3} & 0 \\
							0 & 1-\frac{\gamma_1}{\gamma_3} & \frac{\gamma_1}{\gamma_3} \\
							1 & 0 & 0 
						\end{array}
						\right) \\ 
						A^{(1,2)}_8 & = \left(
						\begin{array}{ccc}
							0 & 1 & 0 \\
							1-\frac{\gamma_1}{\gamma_3} & 0 & \frac{\gamma_1}{\gamma_3} \\
							1 & 0 & 0 \\
						\end{array}
						\right) &
						A^{(1,2)}_9 & = \left(
						\begin{array}{ccc}
							1-\frac{\gamma_1}{\gamma_3} & 0 & \frac{\gamma_1}{\gamma_3} \\
							0 & 1 & 0 \\
							1 & 0 & 0 \\
						\end{array}
						\right) & 
						A^{(1,2)}_{10} & = \left(
						\begin{array}{ccc}
							0 & 1-\frac{\gamma_1}{\gamma_3} & \frac{\gamma_1}{\gamma_3} \\
							1-\frac{\gamma_1}{\gamma_3} & \frac{\gamma_1}{\gamma_3} & 0 \\
							1 & 0 & 0 \\
						\end{array}
						\right).
				\end{aligned}}
			\end{equation}
			All operations therefore interpolate between a permutation at $\gamma_e = \gamma_g$ and a block-upper-triangular matrix when $\gamma_g = 1-\gamma_e = 1$.
			
			We finish consideration of extreme operation flow in $n=3$ by putting forward an independent reproduction of extreme operation structure presented in \cite{mazurek2018decomposability}. Below we skip the labeling, which can be found therein.
			
			\begin{figure}[h]
				\centering
				\includegraphics[width=.5\linewidth]{111_structure.pdf}
				\caption{Flowchart for extreme points of $\mathcal{S}^{\gamma(\beta)}$ for $n=3$ with $\mu = (1,1,1)$, i.e. the nondegenerate three-level case. The structure matches the flowchart given in~\cite{mazurek2018decomposability}.}
				\label{fig:flowchart_111}
			\end{figure}
			
			% \break
			
			For $n=4$ the situation is far more complicated, as generically one will encounter a maximum of 7 different critical temperatures, with merging of distinct temperatures subject not only to degeneracy of energy levels but also specific structure of the energy gaps.  
			
			% \vfill
			
			\begin{table}[ht]
				\centering
				\begin{tabular}{r|c|c|c|c}
					$\mu$ & $\beta \rightarrow \infty$ & $\beta > \beta_{\rm low}$ & $\beta < \beta_{\rm high}$ & $\beta = 0$   \\
					\hline
					$(1,1,1,1)$ & $24$ & $118$ & $262$ & \multirow{8}{*}{$24$} \\
					$(1,1,2)$ & $24$ & $68$ & $200$ &\\
					$(1,2,1)$ & $28$ & $98$ & $200$ & \\
					$(1,3)$ & $24$ & $34$ & $112$ & \\
					$(2,1,1)$ &  $28$ & $128$ & $200$ &\\
					$(2,2)$ & $28$ & $68$ & $124$ &\\
					$(3,1)$ & $34$ & $96$ & $112$ &\\
					$(4)$ & $24$ & $24$ & $24$ &
				\end{tabular}
				\caption{Number of extreme points of the set of Gibbs-preserving operations $\mathcal{S}^{\gamma(\beta)}$ for $n = 4$, depending on the underlying degeneracy structure $\mu$. At infinite temperature, $\beta = 0$, the number is constant and equal to $n!$, and in the zero-temperature limit, $\beta \rightarrow \infty$, it is given by Proposition \ref{prop:extreme_pts}. The resulting values in these two asymptotic regimes, computed explicitly using the Wolfram script available at \cite{github}, exhibit substantial variation with the underlying degeneracy structure.}
				\label{tab:num_points}
			\end{table}
			
			\section{Proofs}
			\label{sec:proofs}
			
			
			\begin{proof}[Proof of Theorem~\ref{thm:general}]
				Since $
				\operatorname*{Li}\mathcal{S}^{\gamma(\beta)}
				\subseteq
				\operatorname*{Ls}\mathcal{S}^{\gamma(\beta)}
				$,
				to show that
				$$
				\operatorname*{Lim}\mathcal{S}^{\gamma(\beta)}
				=
				\mathcal T^{q^1,\ldots,q^k}
				$$
				it suffices to prove the two inclusions
				\begin{equation}
					\label{eq:2_inclusions}
					\operatorname*{Ls}\mathcal{S}^{\gamma(\beta)}
					\subseteq
					\mathcal T^{q^1,\ldots,q^k},
					\qquad
					\mathcal T^{q^1,\ldots,q^k}
					\subseteq
					\operatorname*{Li}\mathcal{S}^{\gamma(\beta)}.
				\end{equation}
				We begin with the upper limit. Suppose that
				$$
				P\in\operatorname*{Ls}\mathcal{S}^{\gamma(\beta)}.
				$$
				Then, there exists a sequence $\beta_\ell\to\infty$ and matrices
				$P^{\beta_\ell}\in S^{\gamma(\beta_\ell)}$
				such that $P^{\beta_\ell}\to P$.
				Since
				$$
				\gamma(\beta_\ell)P^{\beta_\ell}
				=
				\gamma(\beta_\ell),
				$$
				for every block index $b=1,\ldots,k$,
				\begin{equation}
					\label{eq:block-balance}
					\sum_{a=1}^k
					\gamma^{a}(\beta_\ell)
					\bigl(P^{\beta_\ell}\bigr)^{ab}
					=
					\gamma^{b}(\beta_\ell),
				\end{equation}
				where $\gamma^{a}(\beta_\ell)$ and 
				$\bigl(P^{\beta_\ell}\bigr)^{ab}$ are conformable for row-by-column multiplication.
				We first show that all strictly lower-triangular blocks of $P^{\beta_{\ell}}$ vanish in the limit.
				
				Fix indices
				$$
				1\le b<a\le k.
				$$
				Since all matrix entries are nonnegative, equation~\eqref{eq:block-balance} implies
				$$
				\gamma^{a}(\beta_\ell)
				\bigl(P^{\beta_\ell}\bigr)^{ab}
				\le
				\gamma^{b}(\beta_\ell),
				$$
				componentwise. Recall that for $b<a$, $$\|\gamma^b(\beta_{\ell})\|_1\ll \|\gamma^a(\beta_{\ell})\|_1,$$ as $\ell\to\infty$. Therefore,
				$$
				\bigl(P^{\beta_\ell}\bigr)^{ab}\to0,
				$$
				and passing to the limit gives
				$$
				P^{ab}=0,
				\quad \text{for $b<a$}.
				$$
				Thus $P$ is block upper triangular.
				
				We now consider the diagonal blocks. Again from~\eqref{eq:block-balance},
				taking $b=a$, we obtain
				$$
				\gamma^{a}(\beta_\ell)
				\bigl(P^{\beta_\ell}\bigr)^{aa}
				\le
				\gamma^{a}(\beta_\ell).
				$$
				Therefore every diagonal block $\bigl(P^{\beta_\ell}\bigr)^{aa}$
				is $\gamma^a(\beta_{\ell})$-subpreserving. Passing to the limit shows that each diagonal block $P^{aa}$ is $q^a$-subpreserving. We conclude that $
				P\in\mathcal T^{q^1,\ldots,q^k}$.
				This proves
				$$
				\operatorname*{Ls}\mathcal{S}^{\gamma(\beta)}
				\subseteq
				\mathcal T^{q^1,\ldots,q^k}.
				$$
				
				We now prove the other inclusion.
				Let $P\in\mathcal T^{q^1,\ldots,q^k}$. 
				We construct matrices
				$P^{\beta}\in \mathcal{S}^{\gamma(\beta)}$
				such that
				$P^{\beta}\to P$.
				The idea is the following: we keep all rows of $P$ except the last one, and modify the final row so as to enforce exact Gibbs invariance.
				Let
				$$
				p=(P_{n1},\ldots,P_{nn})
				$$
				be the last row of $P$, and let
				$$
				p^{(\beta)}=(p^{(\beta)}_1,\ldots,p^{(\beta)}_n)
				$$
				be defined by
				\begin{equation}
					\label{eq:pbeta_def}
					p^{(\beta)}=\frac{1}{\gamma_n(\beta)}\left[\gamma(\beta)-\gamma(\beta)P\right]+p.
				\end{equation}
				In coordinates,
				\begin{equation}
					\label{eq:pbeta}
					p^{(\beta)}_j
					=
					\frac{
						\gamma_j(\beta)
						-
						\sum_{i=1}^{n-1}
						\gamma_i(\beta)P_{ij}
					}
					{\gamma_n(\beta)},
					\qquad
					j=1,\ldots,n.
				\end{equation}
				Define $P^{\beta}$ by replacing the last row $p$ of $P$ with $p^{(\beta)}$:
				\begin{equation}
					P^{\beta}_{ij}
					=
					\begin{cases}
						P_{ij}, & i<n,\\
						p^{(\beta)}_j, & i=n.
					\end{cases}
				\end{equation}
				By construction, $P^{\beta}$ satisfies
				$$
				\gamma(\beta)P^{\beta}
				=
				\gamma(\beta),
				$$
				for all $\beta\geq0$. 
				We claim that $P^{\beta}$ has row sums equal to $1$, and nonnegative entries for all sufficiently large $\beta$. The claim is true for the first $n-1$ rows of $P^{\beta}$. Since $P$ is row-stochastic, from~\eqref{eq:pbeta_def} we have that the last row $p^{(\beta)}$ of $P^{\beta}$ sums to $1$. 
				Positivity for sufficiently large $\beta$ follows from the $q^a$-subpreserving property of the diagonal blocks $P^{aa}$.
				Hence, we exhibited a family $P^{(\beta)}$ of  $\gamma(\beta)$-preserving row-stochastic matrices such that  $P^{(\beta)}\to P$.
				Thus
				$$
				P\in
				\operatorname*{Li}\mathcal{S}^{\gamma(\beta)}.
				$$
				Since $P\in\mathcal T^{q^1,\ldots,q^k}$ was arbitrary,
				$$
				\mathcal T^{q^1,\ldots,q^k}
				\subseteq
				\operatorname*{Li}\mathcal{S}^{\gamma(\beta)}. 
				$$
				Combining the two inclusions~\eqref{eq:2_inclusions}
				we conclude that
				\begin{equation}
					\operatorname{Lim}S^{\gamma(\beta)}= \mathcal T^{q^1,\ldots,q^k}.
				\end{equation}
				
				
			\end{proof}
			
			
			
			\begin{proof}[Proof of Proposition~\ref{prop:extreme_pts}] The extreme points of $\mathcal{T}_{\mu}$ are $0$--$1$ matrices:
				\begin{equation}
					\operatorname{ext}\left(\mathcal{T}_{\mu}\right)=\left\{P\in\{0,1\}^{n\times n}\colon P\in \mathcal{T}_{\mu}\right\},
				\end{equation}
				as can be shown by contradiction.
				The block-upper-triangular $0$--$1$ matrices $P$  in $\mathcal{T}_{\mu}$ are exactly those that satisfy the following two conditions:
				\begin{enumerate}
					\item each row contains exactly one $1$ (row-stochasticity of $P$);
					\item for every $a=1,\ldots,k$, each column of the diagonal block $P_{aa}$ contains at most one $1$ (column-substochasticity of $P^{aa}$).
				\end{enumerate}
				
				Using the above characterisation, we now  derive the counting formula
				\begin{equation}
					\label{eq:claimed_formula}
					\#\operatorname{ext}\left(\mathcal{T}_{\mu}\right)=\prod_{a=1}^k
					\sum_{i+i'=m_a}
					i!\binom{m_a}{i}^2
					\left(m_{a+1}+\cdots+m_k\right)^{i'}.
				\end{equation}
				The Laguerre representation~\eqref{eq:Laguerre_repr} then follows immediately.
				
				Since the blocks are independent, it suffices to count, for each fixed $a$, the number of admissible matrices of size
				$$
				m_a\times (m_a+\cdots+m_k).
				$$
				
				We distinguish two types of rows.
				
				Suppose that exactly $i$ rows place their $1$ inside the diagonal block $P^{aa}$. The remaining $
				i'=m_a-i$
				rows place their $1$ strictly to the right of the diagonal block.
				There are
				$$
				\binom{m_a}{i}
				$$
				ways to choose the $i$ rows contributing to the diagonal block. Once these rows are selected, their $1$'s must occupy distinct columns among the first $m_a$ columns. Hence the number of possibilities is
				$$
				m_a(m_a-1)\cdots(m_a-i+1)
				=
				i!\binom{m_a}{i}.
				$$
				The remaining $i'$ rows may place their $1$ arbitrarily among the last
				$m_{a+1}+\cdots+m_k$
				columns, giving
				$$
				\left(m_{a+1}+\cdots+m_k\right)^{i'}
				$$
				possibilities.
				
				Therefore the total number of admissible matrices associated with the $a$-th block row is
				$$
				\sum_{i+i'=m_a}
				\binom{m_a}{i}
				\,i!\binom{m_a}{i}\,
				\left(m_{a+1}+\cdots+m_k\right)^{i'}=\sum_{i+i'=m_a}
				i!\binom{m_a}{i}^2
				\left(m_{a+1}+\cdots+m_k\right)^{i'}
				$$
				Multiplying over $a=1,\ldots,k$
				yields~\eqref{eq:claimed_formula}.
			\end{proof}
			
			
			
			
			\begin{proof}[Proof of Theorem~\ref{prop:monotones}] 
				
				Suppose first that 
				$x=yP$ 
				for some $P\in\mathcal T_\mu$.
				Write $P=(P^{ab})_{a,b=1}^k$ according to the block decomposition induced by
				$\mu=(m_1,\ldots,m_k)$.
				Since $P$ is block upper triangular,
				$$
				x^a
				=
				y^aP^{aa}
				+
				z^a,
				\qquad
				a=1,\ldots,k,
				$$
				where
				$$
				z^a:=\sum_{b<a}y^bP^{ba}.
				$$. 
				Since the diagonal block $P^{aa}$ is column-substochastic, weak majorisation theory implies
				$$
				\sum_{i=1}^r
				\bigl(y^aP^{aa}\bigr)_i^\downarrow
				\le
				\sum_{i=1}^r
				(y^a)_i^\downarrow,\quad\text{for all $r=1,\ldots,m_a$.}
				$$
				Moreover, since $z^a$ has nonnegative entries,
				$$
				\sum_{i=1}^r
				(x^a)_i^\downarrow
				=
				\sum_{i=1}^r
				\bigl(y^aP^{aa}+z^a\bigr)_i^\downarrow
				\le
				\sum_{i=1}^r
				\bigl(y^aP^{aa}\bigr)_i^\downarrow
				+
				\|z^a\|_1,
				$$
				because the sum of the $r$ largest entries of a nonnegative vector increases
				by at most the total mass that is added.
				Combining the two inequalities yields
				\begin{equation}
					\label{eq:block}
					\sum_{i=1}^r
					(x^a)_i^\downarrow
					\le
					\sum_{i=1}^r
					(y^a)_i^\downarrow
					+
					\|z^a\|_1,
					\qquad
					r=1,\ldots,m_a.
				\end{equation}
				
				Now let
				$$
				\ell
				=
				m_1+\cdots+m_{a-1}+r,
				\qquad
				1\le r\le m_a.
				$$
				Then
				$$
				\sum_{i=1}^{\ell}x_i^{\mu,\downarrow}
				=
				\sum_{c<a}\|x^c\|_1
				+
				\sum_{i=1}^r(x^a)_i^\downarrow.
				$$
				Applying~\eqref{eq:block},
				\begin{equation}
					\label{eq:intermediate_transf}
					\sum_{i=1}^{\ell}x_i^{\mu,\downarrow}
					\le
					\sum_{c<a}\|x^c\|_1
					+
					\sum_{i=1}^r(y^a)_i^\downarrow
					+
					\|z^a\|_1. 
				\end{equation}
				It therefore remains to compare the total masses of the preceding blocks.
				For every block $c$,
				$$
				\|x^c\|_1
				=
				\|y^cP^{cc}\|_1+\|z^c\|_1.
				$$
				Since $P$ is row-stochastic,
				$$
				\|y^cP^{cc}\|_1
				=
				\|y^c\|_1
				-
				\sum_{d>c}\|y^cP^{cd}\|_1.
				$$
				Hence
				$$
				\|x^c\|_1
				=
				\|y^c\|_1
				-
				\sum_{d>c}\|y^cP^{cd}\|_1
				+
				\|z^c\|_1.
				$$
				Summing over $c<a$ gives
				\begin{equation}
					\label{eq:mass_transf}
					\sum_{c<a}\|x^c\|_1
					=
					\sum_{c<a}\|y^c\|_1
					-
					\sum_{c<a}\sum_{d>c}\|y^cP^{cd}\|_1
					+
					\sum_{c<a}\|z^c\|_1. 
				\end{equation}
				Observe that
				$$
				\|z^c\|_1
				=
				\sum_{b<c}\|y^bP^{bc}\|_1,
				$$
				since all matrices and vectors are nonnegative.
				Hence
				$$
				\sum_{c<a}\|z^c\|_1
				=
				\sum_{b<c<a}\|y^bP^{bc}\|_1.
				$$
				On the other hand,
				$$
				\sum_{c<a}\sum_{d>c}\|y^cP^{cd}\|
				=
				\sum_{c<d<a}\|y^cP^{cd}\|_1
				+
				\|z^a\|_1
				+
				\sum_{\substack{c<a\\ d>a}}\|y^cP^{cd}\|_1.
				$$
				The first term on the right-hand side is precisely
				$\sum_{c<a}\|z^c\|_1$, and it accounts for all transport
				between the first $a-1$ blocks.
				
				
				Substituting this identity into~\eqref{eq:mass_transf} yields
				$$
				\sum_{c<a}\|x^c\|_1
				+\|z^a\|_1
				=
				\sum_{c<a}\|y^c\|_1
				-
				\sum_{\substack{c<a\\ d>a}}
				\|y^cP^{cd}\|_1.
				$$
				Since every summand on the right-hand side is nonnegative,
				$$
				\sum_{c<a}\|x^c\|_1
				+\|z^a\|_1
				\le
				\sum_{c<a}\|y^c\|_1.
				$$
				Therefore, the inequality~\eqref{eq:intermediate_transf} becomes
				$$
				\sum_{i=1}^{\ell}x_i^{\mu,\downarrow}
				\le
				\sum_{c<a}\|y^c\|_1
				+
				\sum_{i=1}^r(y^a)_i^\downarrow
				=
				\sum_{i=1}^{\ell}y_i^{\mu,\downarrow}.
				$$
				
				Since $\ell$ was arbitrary, we conclude that
				$$
				\sum_{i=1}^{\ell}x_i^{\mu,\downarrow}
				\le
				\sum_{i=1}^{\ell}y_i^{\mu,\downarrow},
				\qquad
				\ell=1,\ldots,n.
				$$
				
				Finally, $P$ is row-stochastic, and therefore
				$$
				\sum_{i=1}^n x_i
				=
				\sum_{i=1}^n y_i.
				$$
				Thus equality holds for $\ell=n$, completing the proof.
				
				
				Conversely, assume that
				\begin{equation}\label{eq:InequalitiesProof}
					\sum_{i=1}^\ell x_i^{\mu,\downarrow}
					\le
					\sum_{i=1}^\ell y_i^{\mu,\downarrow},
					\qquad
					\ell=1,\ldots,n,
				\end{equation}
				with equality at $\ell=n$. 
				We construct $P\in\mathcal T_\mu$ such that $x=yP$.
				
				Proceed block by block. Recalling the notation \eqref{eq:NotationMuOrdering}, the inequalities \eqref{eq:InequalitiesProof} can be written as 
				\begin{equation}\label{eq:InequalitiesFirstBlock}
					\sum_{j=1}^{\ell}(x^{1})^{\downarrow}_j\leq \sum_{j=1}^{\ell}(y^1)^{\downarrow}_j, \qquad \ell=1,\dots,m_1
				\end{equation}
				for $a=1$ and 
				\begin{equation}\label{eq:generalInequality}
					\|x^{1}\|_{1}+\dots +\|x^{a-1}\|_{1}+\sum_{j=1}^{\ell} (x^a)^{\downarrow}_{j} \leq  \|y^{1}\|_{1}+\dots +\|y^{a-1}\|_{1}+\sum_{j=1}^{\ell} (y^a)^{\downarrow}_{j} , \qquad \ell=1,\dots,m_a
				\end{equation}
				for $1<a\leq k$. The inequalities \eqref{eq:InequalitiesFirstBlock} imply the existence of a column-substochastic matrix $P^{11}$ such that $x^1=y^1 P^{11}$ by Hardy--Littlewood--Pólya theorem~\cite{marshall1979inequalities}. For $a>1$, the diagonal blocks $P^{aa}$ can be again obtained analogously, with the additional caution of using the strictly upper-block entries of $P$ to transport the deficit of mass $\sum_{b<a}(\|y^{b}\|_1-\|x^{b}\|_1)$ from the preceding blocks, which is non-negative thanks to the cumulative inequalities. More explicitly, the inequalities \eqref{eq:generalInequality} can be rewritten as 
				\begin{equation}
					\sum_{j=1}^{\ell} (x^a)^{\downarrow}_{j} \leq  \sum_{j=1}^{\ell} (y^a)^{\downarrow}_{j}+\sum_{b<a}(\|y^{b}\|_1-\|x^{b}\|_1) , \qquad \ell=1,\dots,m_a,
				\end{equation}
				then Hardy--Littlewood--Pólya theorem guarantees the existence of a column-substochastic matrix $P^{aa}$ such that
				\begin{equation}\label{eq:additionalTerm}
					x^{a}=y^a P^{aa}+\sum_{b<a}(\|y^{b}\|_1-\|x^{b}\|_1)v^aP^{aa},
				\end{equation}
				with $v^a=(1,0,\dots,0)$. Then, using the fact that for $b<a$ we have found $x^b=\sum_{c<b}y^{c}P^{cb}$, the last term of \eqref{eq:additionalTerm} contains only linear combinations of the components of $y^{1},\dots,y^{a-1}$, which allows the extraction of the blocks $P^{1a},\dots,P^{a-1,a}$. The resulting matrix is block upper triangular, and has column-substochastic diagonal blocks by construction, and one can check that it is also row-stochastic and has non-negative elements. 
			\end{proof}
			
			\begin{proof}[Proof of Proposition~\ref{prop:fate_monotones}]
				Assume that there exists $\beta_0$, such that $
				x\prec^{\gamma(\beta)} y$ for all $\beta>\beta_0$. Then, if $\beta>\beta_0$,
				\begin{equation}
					F_{\beta}(x)=U(x)-\frac1\beta S(x)
					\le
					U(y)-\frac1\beta S(y)=F_{\beta}(y).
				\end{equation}
				Hence, either $
				U(x)<U(y)$,
				or
				$U(x)=U(y) $ but 
				$S(x)\ge S(y)$.
				
				Consider now the inequalities
				$$
				P_\alpha(x)\le P_\alpha(y)
				$$
				for all $\beta>\beta_0$. 
				
				The inequalities read explicitly,
				\begin{equation}
					\sum_j x_j^\alpha e^{-\beta(1-\alpha)E_j}\leq  \sum_j y_j^\alpha e^{-\beta(1-\alpha)E_j}
				\end{equation}
				If $1-\alpha>0$, the dominant contribution to $P_{\alpha}$ comes from the lowest energy sectors, yielding the reverse lexicographic condition~\eqref{eq:revlex-condition}. 
				If $1-\alpha<0$, the highest energy sectors dominate, and we have the lexicographic condition~\eqref{eq:lex-condition}.
				
			\end{proof}
			
			\bibliographystyle{alpham}
			\bibliography{References}
			
		\end{document}
